

\section{Thesis proposal}

\subsection{Introduction}

Quantum computers promise vast computational power, made possible by the surprising laws of quantum mechanics. The first quantum computers will certainly be extremely expensive. Consequently, their services will be available online, as is already the case for several popular quantum platforms like IBM. However, a company owning a secret high-value algorithm will undoubtedly try to protect it against a dubious cloud provider: it is thus of utmost importance to \gras{establish trust between users and quantum providers} by answering the following questions:

\begin{center}
\vspace{-.5\baselineskip}
  \emph{Can a user delegate a computation to a remote server while making sure that their private data/algorithms stay hidden to a malicious server and that the obtained result is correct?}
  \vspace{-.5\baselineskip}
\end{center}

The \gras{goal of this thesis will be to answer this fundamental question}\===overcoming severe limitations in the existing literature as explained below\===and \gras{explore other cryptographic applications} that are made possible thanks to the unique properties of quantum mechanics. To this end, we structure the thesis schedule around \gras{three one-year projects} described below.

\subsection{Project I: Verifiable blind quantum computing for non-deterministic computations}

The first project aims to answer the above question by providing the first protocol that can blindly (i.e.\ by hiding to the server the inputs, outputs, and algorithm of the client) delegate a circuit execution to a remote server while allowing the classical (i.e.\ non-quantum) client to verify the result provided by the server. While this fundamental problem was already extensively studied~\cite{CCKW19_QFactoryClassicallyInstructedRemote,mahadev_classical_2018}, the existing literature has some very severe limitations since none of these results provides both blindness and verifiability with a classical client. The most promising candidate is~\cite{GV19_ComputationallySecureComposableRemote} (FOCS 2019), but it does not only have poor (polynomial) security guarantees, but \gras{we have also found with our candidate Louis two fundamental mistakes} in this work, invalidating its conclusions. Given the importance of this result, we aim to fix these issues during the first part of this thesis, based on very promising results obtained with Louis during his Master 1 internship.

%that pretends to have both a verifiable and blind protocol with classical clients, but this work has two major drawbacks: first, the security scales only polynomially with the number of rounds (we can attack the protocol with probability $\approx\frac{1}{\sqrt{\text{nb rounds}}}$) while we typically want exponential security. But more importantly, we also believe that \gras{this article has some fundamental mistakes} invalidating their construction in a way that cannot easily be fixed. Given the importance of this result and its wide adoption in the community, we aim to fix these issues during the first part of this thesis.

% Our candidate, Louis \textsc{Triouleyre Roberjot}, did his Master 1 internship with Léo \textsc{Colisson Palais} (supervisor of this thesis proposal) on this exact question and already has really promising early results. As such, we are confident that Louis is an excellent candidate and that he will be able to efficiently tackle the above problem. % already said below

\subsection{Project II: One-Time Programs}

One-Time Programs (OTPs)~\cite{BGS13_QuantumOneTimePrograms} allow an adversary to execute at most once a given program, paving the way to many applications like one-time signatures. OTPs are impossible to realize without hardware assumptions, but the assumptions used nowadays are very limiting since we typically need to send a physically impenetrable device. While \cite{gupte_quantum_2025} managed to avoid hardware assumptions but their results only work for a sub-class of randomized circuits, \cite{Liu23_DepthBoundedQuantumCryptography} is very generic but assumes that there exists a short time after which any quantum state will decohere. Instead, in this project we want to \gras{build generic OTPs using a weaker hardware assumption}, namely that there exists some functions that can be classically computed but that cannot, in practice, be run in quantum superposition, leading to practical OTPs in the noisy intermediate-scale quantum (NISQ) era.

%In a first step we will therefore obtain such a protocol by relying on a notion of ``classically queriable random oracle'' (cqRO) and combine it with methods coming from error correcting codes. This part of the project has \gras{low risk} as we have promising candidates of protocols. It is thus \gras{well-suited for a master's student}. In a more exploratory approach, we will continue this WP by studying methods to obtain quantum-memory-hard functions (following the idea of memory-hard classical functions which are functions that require a large amount of memory to run), i.e.\ functions that require a large quantum memory to be executed in superposition, making them impossible to implement on a quantum device. In a more exploratory direction, we also want to explore if methods introduced in \cite{gupte_quantum_2025} combine with methods we introduced in \cite{CMS23_ObliviousTransferZeroKnowledge} can help us to obtain quantum money.

% OTPs have many of applications, for instance to allow a user to use a software only once (e.g.\ to try it) or in the context of signature delegation. By making OTPs finally accessible in the short term (NISQ era), we will provide some applications to quantum computing (beyond Quantum Key Distribution) that can already be implemented now, without waiting for the existence of full-fledged quantum computers.

% Moreover, we will find the \gras{first classically verifiable QPV protocols with security guarantees that apply to realistic settings}. As these protocols are especially promising in the near term, we will advertise our results to experimental groups that could be interested in realizing them in a future experiment, which would mark an important step towards practical deployment of QPV.

\subsection{Project III: Classically-verifiable quantum position verification}

Quantum Position Verification (QPV) allows a party to prove to a set of verifiers that they are in a given physical location, using notably bounds on the speed of light. CTWhile a fully classical protocol is always subject to attacks\===hence the need for quantum protocols\===little work has been done to minimize the quantum requirements, as most existing protocols require all parties to have quantum capabilities and to quantumly communicate (which is extremely impractical). In this project, we aim to improve over~\cite{LLQ21_BeatingClassicalImpossibility} to obtain the first classically verifiable QPV protocol, i.e., involving only classical communication, that is secure under realistic conditions (plain model) against attackers with polynomial entanglement, paving the way to practical QPV. %To attain this aim, we will need easier PoQs with more realistic security guarantees, since any classically verifiable QPV protocols needs to rely on PoQs \cite{LLQ21_BeatingClassicalImpossibility}. Therefore, we will investigate how to use both constructions from the QFactory protocol for RSP that PI Léo Colisson Palais has worked on
%\cite{CCKW19_QFactoryClassicallyInstructedRemote} and the security analysis for PoQ protocols from WP 1.1. This part of the project has \gras{medium risk} due to its ambition, which is mitigated by the expertise of PIs Léo Colisson Palais and Andreas Bluhm in both PoQs \cite{CCKW19_QFactoryClassicallyInstructedRemote} and QPV \cite{bluhm2022single, allerstorfer2023making, colisson_palais_quantum_2025, bluhm2025complexity}.

\section{Connection with the PERSYVAL LabEx}

Our project falls within the Persyval's MFA (\emph{Mathématiques : des fondements aux applications}) axis, since this PhD thesis focuses on \emph{quantum cryptography}, which is the first domain \mylink{https://persyval.univ-grenoble-alpes.fr/recherche/axes/mfa}{highlighted in the MFA axis}. Advanced mathematical tools are indeed involved in every aspect of the proposal: for instance, the VFactory project will rely on the fundamental structure of Hilbert spaces\===described notably by the Jordan's lemma that explains how two projectors on Hilbert spaces can be decomposed into direct sums on one or two-dimensional subspaces\===and lattice-based problems.

\section{Complementarity between the supervisors}

The three supervisors, Léo \textsc{Colisson Palais}, Mehdi \textsc{Mhalla} and Andreas \textsc{Bluhm}, have a very complementary expertise that will turn out to be particularly helpful in the context of this thesis: L.~\textsc{Colisson Palais}'s expertise focuses on methods involving computational properties (LWE, random oracle model…), while M.\ \textsc{Mhalla} and A.\ \textsc{Bluhm} have a very solid expertise on quantum information theory. The three projects uniquely mix these computational and statistical properties: L.\ \textsc{Colisson Palais} is the author of \cite{CCKW19_QFactoryClassicallyInstructedRemote} (project I's starting point), and the rest of the project will involve QRAC and graph-states, extensively studied by M.\ \textsc{Mhalla} and A.\ \textsc{Bluhm}. Additionally, L.\ \textsc{Colisson Palais} and M.\ \textsc{Mhalla} already initiated a projects on One-Time Program (project II) involving the random oracle, graph-states and error-correcting codes. Finally, QPV (last project) is the main research area of A.\ \textsc{Bluhm}, and we will incorporate technics based on LWE.

Lastly, both Léo and Mehdi supervised Louis (resp.\ in Master 1 and Master 2) and can assert that Louis is an excellent and extremely motivated student that will be perfectly suited for this proposal as he already has promising results on project I.

\section{References}
\printbibliography[heading=none]


\end{document}
